\documentclass[11pt]{article}

\usepackage{amsmath}
\usepackage{amssymb}
\usepackage{graphicx}
\usepackage{booktabs}
\usepackage{float}
\usepackage{geometry}
\usepackage{caption}
\usepackage{enumitem}

\geometry{margin=1in}

\title{OverclusterDRG: A Deterministic, LP-Free Heuristic\\
for Robust $k$-Center with Outliers}
\author{}
\date{}

\begin{document}

\maketitle

%% ================================================================
\section{Problem Definition}
%% ================================================================

Let:
\begin{itemize}
    \item $X = \{x_1, \dots, x_N\}$ be the dataset of $N$ points,
    \item $k$ be the number of centers to place,
    \item $z$ be the maximum number of outliers allowed.
\end{itemize}

For any center set $C$, the distance from a point $x$ to its nearest center is
\[
d(x, C) = \min_{c \in C}\, d(x, c).
\]

If $S \subseteq X$ with $|S| \le z$ is the outlier set, the clustering radius is
\[
R(C, S) = \max_{x \in X \setminus S}\, d(x, C).
\]

The objective is
\[
\min_{\substack{C \subseteq X,\; |C|=k \\ S \subseteq X,\; |S| \le z}} R(C,S).
\]

We call the optimal value $\mathrm{OPT}$.

For a fixed center set $C$, the optimal outlier set is always the $z$ farthest points. So we define
\[
r_z(C) \;=\; \text{the }(z{+}1)\text{-th largest value in } \{d(x,C) : x \in X\},
\]
which is the radius after discarding the $z$ worst points. Then $\mathrm{OPT} = \min_{|C|=k}\, r_z(C)$.

\subsection{LP Lower Bound}

We use a linear program as a lower bound baseline. For a fixed candidate radius $r$, define
\[
N(j, r) = \{ i \mid d(i,j) \le r \}.
\]

The LP feasibility formulation is:
\begin{align*}
\text{Find } & x_{ij},\, y_i,\, c_j \\
\text{s.t.}\quad
& \sum_{i \in N(j,r)} x_{ij} \ge c_j \quad \forall j \\
& x_{ij} \le y_i \quad \forall i,j \\
& \sum_i y_i \le k \\
& \sum_j c_j \ge N - z \\
& 0 \le x_{ij},\, y_i,\, c_j \le 1
\end{align*}

We binary-search the sorted list of pairwise distances and solve the LP at each candidate $r$. The smallest feasible $r$ is $R_\mathrm{LP}$, which satisfies $R_\mathrm{LP} \le \mathrm{OPT}$. Every ratio $r_z(C)/R_\mathrm{LP}$ reported in the experiments upper-bounds the true approximation ratio against $\mathrm{OPT}$.

\subsection{Two Optima}

The theory uses two different optimal values:
\begin{itemize}
    \item $\mathrm{OPT}_{k,z}$ = the robust $k$-center optimum (discard up to $z$ outliers). This is what we are solving.
    \item $\mathrm{OPT}_{m,0}$ = the standard $m$-center optimum with no outliers, covering all $N$ points.
\end{itemize}

The relationship between $\mathrm{OPT}_{k+z,0}$ and $\mathrm{OPT}_{k,z}$ is the key step in the analysis.


%% ================================================================
\section{Dataset and Preprocessing}
%% ================================================================

\subsection{Adult Dataset}

The UCI Adult (Census Income) dataset contains 14 attributes and a binary income label. This report uses only the training split and only the 6 numeric attributes:
\[
\texttt{age, fnlwgt, education-num, capital-gain, capital-loss, hours-per-week}.
\]
The income label is not used. Distance is Euclidean on z-scored numeric vectors.

\subsection{Experimental Grid}

\[
N \in \{500, 1000, 2000, 5000, 10000, 20000\},\quad
k \in \{10, 20, 50\},\quad
z/N \in \{0.10, 0.20\}.
\]

Total: $6 \times 3 \times 2 = 36$ configurations. Secondary datasets (Diabetes Indicators, Covertype) use the same $k$ and outlier-rate values with different size grids.


%% ================================================================
\section{Algorithms}
%% ================================================================

\subsection{RKC-Gonzalez}

At each of the $k$ steps, instead of picking the single farthest point, pick the $(z{+}1)$-th farthest, skipping the $z$ candidates most likely to be outliers.

\begin{verbatim}
Input: X (N points), k, z
Output: C (centers), R (radius)

1. c1 = point nearest to centroid of X
   C = {c1}
   dist[x] = d(x, c1) for all x

2. For i = 2 to k:
      Sort X by dist[] descending
      c_i = point at rank z+1   // skip top-z outlier candidates
      C = C + {c_i}
      dist[x] = min(dist[x], d(x, c_i)) for all x

3. Sort {dist[x] : x in X} descending as d_1 >= d_2 >= ... >= d_N
   R = d_{z+1}   // drop z farthest, take next

4. Return C, R
\end{verbatim}

This runs in $O(Nk \log N)$ and has no approximation guarantee.

\subsection{Charikar 3-Approximation}

\begin{verbatim}
Input: X (N points), k, z
Output: C (centers), R_out (certified radius)

1. Collect all pairwise distances; sort them as candidates.
2. Binary-search for the smallest feasible radius r:

   Feasibility check at radius r:
      U = X  (uncovered points)
      C = empty
      While |U| > z and |C| < k:
         For each q in U, count points in B(q, r)   // ball of radius r
         q* = argmax count
         C = C + {q*}
         Remove all points in B(q*, 3r) from U      // 3r coverage
      If |U| <= z: feasible. Else: infeasible.

3. r* = smallest feasible r
4. Return C from step 2, R_out = 3r*
\end{verbatim}

Any point within $r$ of a chosen center $q^*$ is within $3r$ of $q^*$, giving the $3$-approximation guarantee. Runtime scales as $O(N^{2.12})$ empirically.

\subsection{Ding et al. Randomised Algorithm}

\begin{verbatim}
Input: X (N points), k, z
       T = ceil((z+1) * ln(100))   // trials for 99% success
Output: best center set found

C_best = empty;  r_best = infinity

For t = 1 to T:
   S_t = random sample of z+1 points from X
   C_t = Gonzalez(X \ S_t, k)   // standard k-center on reduced set
   r_t = r_z(C_t, X)            // evaluate on full dataset
   If r_t < r_best:
      C_best = C_t;  r_best = r_t

Return C_best
\end{verbatim}

With $T = \lceil(z{+}1)\ln 100\rceil$ trials, the probability that at least one sample $S_t$ intersects the true optimal outlier set is $\ge 99\%$, giving a $2$-approximation with high probability.

\subsection{OverclusterDRG (OC-Forward and OC-Local)}

\begin{verbatim}
Input: X (N points), k, z,  O = z (overcount)
Output: C (k centers)

// Phase 1: build pool of k+O points using standard Gonzalez
q1 = point nearest to centroid of X
Q = {q1};  dist[x] = d(x, q1) for all x
For i = 2 to k+O:
   q_i = argmax_x dist[x]        // farthest point, no skipping
   Q = Q + {q_i}
   dist[x] = min(dist[x], d(x, q_i)) for all x

// Phase 2a: forward greedy selection from pool (= OC-Forward)
C = {q1};  mnd[x] = d(x, q1) for all x
For i = 2 to k:
   q* = pool point in Q\C that minimises r_z when added to C
   C = C + {q*}
   mnd[x] = min(mnd[x], d(x, q*)) for all x

// Phase 2b: 1-swap local search (additional step for OC-Local)
Repeat:
   improved = false
   For each pair (c_i, q) in C x (Q\C), in pool-index order:
      If r_z(C - {c_i} + {q}) < r_z(C):
         C = C - {c_i} + {q}
         improved = true
         Break
Until not improved

Return C
\end{verbatim}

Phase 1 is standard Gonzalez run for $k+z$ steps (no rank-skipping). The extra $z$ steps let the pool absorb the outliers that farthest-point inevitably selects. Phase 2a greedily picks $k$ centers from the pool. Phase 2b polishes with $1$-swap until no improvement is possible. The full OC-Local algorithm completes all three phases.


%% ================================================================
\section{Why OC-Local Works: Theory}
%% ================================================================

This section explains the main theoretical result --- the phase transition at $O = z$ --- in plain terms, with each proof step written out.

\subsection{The Core Idea}

Gonzalez picks the farthest point at every step. Outliers are far from the data by definition, so they attract the farthest-point rule. In a dataset with $z$ outliers, Gonzalez wastes $z$ of its $k + O$ steps selecting them before getting to real cluster representatives.

That leaves $k + O - z$ useful steps. With $O = z$ this is exactly $k$ --- one representative per cluster, none to spare. Run one fewer step ($O = z-1$) and one cluster has no representative in the pool; the algorithm can fail arbitrarily. Run $O \ge z$ steps and the pool is guaranteed to contain a $3 \cdot \mathrm{OPT}$ solution.

\subsection{Lemma 1: What the Pool Contains}

Let $Q = \mathrm{Gonzalez}(X, k+O)$ be the pool. We claim there exists a $k$-subset $C_0 \subseteq Q$ satisfying
\[
r_z(C_0) \;\le\; \mathrm{OPT} + 2 \cdot \mathrm{OPT}_{k+O,0}.
\]

\textbf{Proof.} Let $(C^*, S^*)$ be an optimal solution achieving radius $\mathrm{OPT}$, with $|C^*| = k$ centers and $|S^*| = z$ outliers.

For each optimal center $c^*_j \in C^*$, let $q_j \in Q$ be the nearest pool point. The Gonzalez algorithm is a 2-approximation for standard $k$-center, so at $k+O$ centers:
\[
d(c^*_j,\, q_j) \;\le\; 2 \cdot \mathrm{OPT}_{k+O,0}.
\]

Set $C_0 = \{q_1, \dots, q_k\}$. For any inlier $x \notin S^*$ whose nearest optimal center is $c^*_j$, apply the triangle inequality:
\begin{align*}
d(x, C_0) &\le d(x, q_j) \\
           &\le d(x, c^*_j) + d(c^*_j, q_j) \\
           &\le \mathrm{OPT} + 2 \cdot \mathrm{OPT}_{k+O,0}.
\end{align*}

Discarding the same outlier set $S^*$ gives $r_z(C_0) \le \mathrm{OPT} + 2 \cdot \mathrm{OPT}_{k+O,0}$. \hfill$\square$

\subsection{Lemma 2: Comparing the Two Optima}

For $O \ge z$:
\[
\mathrm{OPT}_{k+O,0} \;\le\; \mathrm{OPT}_{k,z}.
\]

\textbf{Proof.} Let $(C^*, S^*)$ be optimal for the robust problem, achieving radius $\mathrm{OPT} = \mathrm{OPT}_{k,z}$. Construct a $(k{+}z)$-center solution with no outliers:
\[
C' = C^* \cup \{s : s \in S^*\}, \qquad |C'| = k + z.
\]

Check coverage for every point in $X$:
\begin{itemize}
    \item For any inlier $x \notin S^*$: $d(x, C') \le d(x, C^*) \le \mathrm{OPT}$.
    \item For any outlier $s \in S^*$: $d(s, C') = 0$ since $s$ is its own center.
\end{itemize}

All $N$ points are within radius $\mathrm{OPT}$ of $C'$, so
\[
\mathrm{OPT}_{k+z,0} \le \mathrm{OPT}.
\]

Since adding more centers can only reduce the optimum, for $O \ge z$:
\[
\mathrm{OPT}_{k+O,0} \le \mathrm{OPT}_{k+z,0} \le \mathrm{OPT}. \hfill\square
\]

\subsection{Theorem 1: $O \ge z$ Is Sufficient}

For $O \ge z$, the pool contains a $k$-subset with $r_z \le 3 \cdot \mathrm{OPT}$.

\textbf{Proof.} From Lemma 1:
\[
r_z(C_0) \le \mathrm{OPT} + 2 \cdot \mathrm{OPT}_{k+O,0}.
\]

From Lemma 2 (since $O \ge z$):
\[
\mathrm{OPT}_{k+O,0} \le \mathrm{OPT}.
\]

Substituting the second inequality into the first:
\[
r_z(C_0) \le \mathrm{OPT} + 2 \cdot \mathrm{OPT} = 3 \cdot \mathrm{OPT}. \hfill\square
\]

\subsection{Theorem 2: $O < z$ Is Not Enough}

For any $k \ge 2$, $z \ge 1$, and $M > 0$, there exists an instance with $\mathrm{OPT} = \varepsilon$ on which the pool with $O = z - 1$ gives
\[
\min_{C \subseteq Q,\,|C|=k} r_z(C) \ge M\varepsilon.
\]

\textbf{Construction.} Set $D = 2(M{+}1)\varepsilon$.
\begin{itemize}
    \item Place $k$ clusters, each with $m > z$ inliers packed within radius $\varepsilon$ of cluster centers $c^*_1, \dots, c^*_k$ spread $D$ apart along one axis.
    \item Place $z$ outliers at distance $R \gg D$ along orthogonal axes.
\end{itemize}

The optimal solution covers all $k$ clusters and discards the $z$ outliers, achieving $\mathrm{OPT} = \varepsilon$.

Since each outlier is at distance $R \gg D$ from any cluster inlier, Gonzalez selects all $z$ outliers before any second cluster gets a representative. With $O = z-1$, the pool contains: one starting cluster inlier, $z$ outliers, and only $k-2$ cluster representatives --- leaving one cluster with no pool point.

Now consider any $k$-subset $C \subseteq Q$. There are two cases.

\textit{Case 1: $C$ contains an outlier $o_i$.} Then $o_i$ occupies one of the $k$ center slots. The remaining $k-1$ centers must cover $k$ clusters, so one cluster is still left with no nearby center. That cluster has $m > z$ inliers, so they cannot all be discarded within the budget of $z$. The radius is at least $D - 2\varepsilon$.

\textit{Case 2: $C$ contains no outlier.} Then $C$ draws all $k$ centers from the $k-1$ cluster representatives in the pool. One cluster has no representative, so its inliers must be covered by the nearest represented cluster center. The two cluster centers are at least $D$ apart; each inlier is within $\varepsilon$ of its own cluster center, so the distance from any inlier of the missing cluster to any pool representative is at least $D - \varepsilon - \varepsilon = D - 2\varepsilon$. Again the cluster has $m > z$ inliers, so not all can be discarded. The radius is at least $D - 2\varepsilon$.

In both cases $r_z(C) \ge D - 2\varepsilon = 2(M+1)\varepsilon - 2\varepsilon = 2M\varepsilon \ge M\varepsilon$. \hfill$\square$

\textbf{Corollary (Phase Transition).} $O = z$ is the exact threshold: below it the approximation ratio is unbounded; at and above it the pool always contains a $3 \cdot \mathrm{OPT}$ subset.


%% ================================================================
\section{Experiments}
%% ================================================================

\subsection{Setup}

Algorithms compared:
\begin{enumerate}
    \item \textbf{LP} (lower bound only --- not an algorithm)
    \item \textbf{Charikar et al.} 3-approximation, reported at achieved radius
    \item \textbf{RKC-Gonzalez} heuristic (picks $(z{+}1)$-th farthest at each step)
    \item \textbf{OC-Forward} (pool + greedy, no local search)
    \item \textbf{OC-Local} (pool + greedy + 1-swap)
    \item \textbf{Ding et al.} randomised, $T = \lceil(z{+}1)\ln 100\rceil$ trials
\end{enumerate}

All LP bounds are solved exactly using Gurobi 12.0.3 (feasibility tolerance $10^{-6}$). Runtimes range from seconds at small $N$ to 94 hours at $N = 20{,}000$. The approximation ratio reported for each algorithm is radius / $R_\mathrm{LP}$; lower is better and $1.000$ is optimal.

\subsection{Approximation Quality on UCI Adult}

\begin{table}[H]
\centering
\begin{tabular}{lcccccc}
\toprule
Algorithm & Mean & Std & Min & Max & Med.\ time & Max time \\
\midrule
LP oracle     & 1.000 & ---   & 1.000 & 1.000 & --- & --- \\
\textbf{OC-Local}  & \textbf{1.147} & 0.047 & 1.088 & 1.265 & 0.36 s & 57 s \\
OC-Forward    & 1.158 & 0.047 & 1.104 & 1.283 & 0.26 s & 52 s \\
Charikar      & 1.173 & 0.056 & 1.104 & 1.306 & 3.60 s & 564 s \\
Ding          & 1.235 & 0.032 & 1.179 & 1.331 & 2.22 s & 431 s \\
RKC-Gonzalez  & 1.280 & 0.049 & 1.197 & 1.370 & 0.002 s & 0.02 s \\
\bottomrule
\end{tabular}
\caption{Approximation ratios (radius / LP) and runtimes on 36 UCI Adult configurations.}
\end{table}

Table 1 ranks the algorithms by quality. RKC-Gonzalez is fastest by far but pays for it --- a mean ratio of $1.280$ versus OC-Local's $1.147$. Charikar gets close to OC-Local in quality ($1.173$) but at a median $10\times$ slowdown; Ding falls between the two and introduces randomness. The $1$-swap step in OC-Local costs a median $0.10$ s and improves quality in every single configuration over OC-Forward, so it is always worth running.

\subsection{Per-Configuration Wins}

Averages can hide losses on individual configurations. Looking at each of the 36 separately:
\begin{itemize}
    \item OC-Local beats Charikar in $33/36$ cases at a median $7\times$ speedup.
    \item OC-Local beats Ding in $35/36$ cases at a median $8\times$ speedup.
    \item OC-Local beats OC-Forward in $36/36$ cases.
\end{itemize}

The three configurations where Charikar wins are all at small $N$ ($\le 1000$) and small $k$, where Charikar's quadratic search is still affordable and occasionally finds a tighter radius.

\subsection{Statistical Testing}

To check whether these gaps are real or could arise from random variation, we run a paired comparison for each baseline.

\paragraph{Setup.} For each configuration $i \in \{1, \dots, 36\}$, compute the paired difference:
\[
d_i = \frac{R_{B,i}}{R_{\mathrm{LP},i}} - \frac{R_{\mathrm{OC},i}}{R_{\mathrm{LP},i}}.
\]
Positive $d_i$ means OC-Local has a lower LP ratio on that configuration --- it is better. We test:
\[
H_0: \mu_d = 0 \quad \text{(no difference)} \qquad \text{vs} \qquad H_1: \mu_d > 0 \quad \text{(OC-Local better)}.
\]

\paragraph{$\bar{d}$ (mean difference).} The average of all 36 paired differences:
\[
\bar{d} = \frac{1}{n}\sum_{i=1}^{n} d_i.
\]
This just describes how much better OC-Local is on average. It says nothing about whether the gap could be noise.

\paragraph{$t$-statistic.} Measures how many standard errors $\bar{d}$ sits above zero:
\[
t = \frac{\bar{d}}{s_d / \sqrt{n}}, \qquad s_d^2 = \frac{1}{n-1}\sum_{i=1}^{n}(d_i - \bar{d})^2.
\]
A large $t$ means the gap is large relative to the variability across configurations.

\paragraph{$p$-value.} The probability of observing a $t$ this large purely by chance, assuming $H_0$ is true:
\[
p = \Pr(T \ge t_{\mathrm{obs}} \mid H_0).
\]
A small $p$ means the gap is very unlikely to be random. We use $\alpha = 0.01$ as the significance threshold.

\paragraph{Wilcoxon signed-rank test.} The $t$-test assumes the differences $d_i$ are roughly normally distributed. The Wilcoxon test makes no such assumption --- it replaces the actual values with their ranks and checks whether positive differences consistently outrank negative ones. Let $R^+$ be the sum of ranks of positive differences and $R^-$ the sum of ranks of negative ones. The test statistic is $W = \min(R^+, R^-)$; a small $W$ (most ranks on the positive side) gives a small $p$-value. The Wilcoxon $p$ reported here is two-sided.

Since all 36 configurations share the same base dataset, the $t$-test overstates independence. The Wilcoxon test is more conservative and handles this better.

\begin{table}[H]
\centering
\begin{tabular}{lcccc}
\toprule
Comparison & $\bar{d}$ & $t$ & $p$ & Wilcoxon $p$ \\
\midrule
OC-Local vs RKC-Gonzalez & $+0.134$ & $+13.3$ & $<10^{-14}$ & $2.9\times10^{-11}$ \\
OC-Local vs Ding         & $+0.089$ & $+14.3$ & $<10^{-15}$ & $5.8\times10^{-11}$ \\
OC-Local vs Charikar     & $+0.026$ & $+5.7$  & $<10^{-5}$  & $3.1\times10^{-8}$ \\
OC-Local vs OC-Forward   & $+0.012$ & $+7.9$  & $<10^{-8}$  & $2.9\times10^{-11}$ \\
\bottomrule
\end{tabular}
\caption{Paired comparisons ($n=36$). Positive $\bar{d}$ means OC-Local is better.
Both tests agree on every row.}
\end{table}

Every gap is significant well below $\alpha = 0.01$ under both tests. The smallest gap --- OC-Local versus its own initialisation OC-Forward --- still gives $p < 10^{-8}$, confirming the $1$-swap step contributes real improvement, not noise.

\subsection{Runtime Scaling}

Runtime was fit using log-log linear regression:
\[
\log T = a + b \log N + \varepsilon,
\]
where $b$ is the empirical scaling exponent. The coefficient of determination is
\[
R^2 = 1 - \frac{\sum_i (y_i - \hat{y}_i)^2}{\sum_i (y_i - \bar{y})^2},
\quad y_i = \log T_i,\quad \hat{y}_i = a + b\log N_i.
\]

\begin{table}[H]
\centering
\begin{tabular}{lccc}
\toprule
Algorithm & Exponent $b$ & $R^2$ \\
\midrule
OC-Local     & 1.85 & 0.976 \\
Charikar     & 2.12 & 0.999 \\
RKC-Gonzalez & 0.80 & 0.939 \\
Ding         & 1.10 & --- \\
\bottomrule
\end{tabular}
\caption{Empirical runtime scaling exponents.}
\end{table}

OC-Local and Charikar are in the same asymptotic class but OC-Local's constant is much smaller. Ding's exponent looks gentle ($N^{1.10}$) but the trial multiplier $(z{+}1)\ln 100$ dominates in practice: at $N=20{,}000$, $k=50$, $z=20\%$, Ding needs 18,423 trials and 431 seconds while OC-Local finishes in 57 seconds.

\subsection{Radius Comparison Averaged over $N$}

\begin{table}[H]
\centering
\begin{tabular}{cccccc}
\toprule
$k$ & Outlier \% & LP & OC-Local & Charikar & Gonzalez \\
\midrule
10 & 10\% & 1.807 & 2.029 & 2.101 & 2.373 \\
10 & 20\% & 1.479 & 1.646 & 1.674 & 1.869 \\
20 & 10\% & 1.467 & 1.710 & 1.736 & 1.904 \\
20 & 20\% & 1.213 & 1.356 & 1.390 & 1.499 \\
50 & 10\% & 1.118 & 1.348 & 1.369 & 1.458 \\
50 & 20\% & 0.915 & 1.053 & 1.082 & 1.159 \\
\bottomrule
\end{tabular}
\caption{Average radius across all $N$ per $(k, \text{outlier \%})$ bucket.
OC-Local is best in every row.}
\end{table}

\subsection{Does $1$-Swap Actually Fail?}

We ran OC-Local and exhaustive pool search on 400 random Adult subsamples ($N \le 1000$, $k \in \{5,10\}$, $z \le 10$).

OC-Local matched exhaustive in 389 of 400 cases (97.2\%). The 11 failures all had near-equal inter-cluster distances --- the exact adversarial structure from the necessity proof. Maximum gap: 4.82\%. This structure essentially never appears in real data.

\subsection{Pool Size: Confirming the Phase Transition}

We swept $O$ from 1 to 2000 across 72 configurations on Diabetes and Covertype.

\begin{table}[H]
\centering
\begin{tabular}{lccc}
\toprule
Pool setting & Within 5\% of best & Mean gap & Max gap \\
\midrule
$O = z$      & 53/72 & $+3.7\%$ & $+11.8\%$ \\
$O = 2z$     & 69/72 & $+1.3\%$ & $+8.2\%$ \\
$O \approx 3.7z$ & 72/72 & ${<}1\%$  & ${<}3\%$ \\
\bottomrule
\end{tabular}
\caption{Quality achieved at different pool sizes across 72 configurations.}
\end{table}

Quality drops steeply below $O = z$ in every configuration, as the theory predicts. Above $O = z$, quality keeps improving modestly; the mean empirical knee is at $1.55z$. The theory says the pool \emph{contains} a good solution from $O = z$ onward; how much $1$-swap \emph{finds} depends on more pool.

The strongest predictor of how far past $z$ quality keeps improving is $z/k$ (Pearson $r = -0.60$, $p < 10^{-9}$). Dimensionality has essentially no effect ($r = +0.04$, $p = 0.72$).

\subsection{Other Datasets}

\begin{table}[H]
\centering
\begin{tabular}{llcccc}
\toprule
Dataset & Algorithm & Mean & Std & Min & Max \\
\midrule
Diabetes       & RKC-Gonzalez & 1.923 & 0.256 & 1.449 & 2.414 \\
($d=8$, $n=36$)& Ding         & 1.227 & 0.048 & 1.143 & 1.330 \\
               & Charikar     & 1.201 & 0.044 & 1.145 & 1.313 \\
               & OC-Local     & \textbf{1.155} & 0.043 & 1.093 & 1.275 \\
\midrule
Covertype      & RKC-Gonzalez & 1.671 & 0.173 & 1.418 & 1.960 \\
($d=10$,$n=30$)& Ding         & 1.213 & 0.049 & 1.121 & 1.300 \\
               & Charikar     & 1.227 & 0.060 & 1.134 & 1.374 \\
               & OC-Local     & \textbf{1.171} & 0.051 & 1.095 & 1.294 \\
\bottomrule
\end{tabular}
\caption{LP ratios on secondary datasets (exact LP bounds).}
\end{table}

OC-Local is the best practical algorithm on both datasets. The ordering is the same as on Adult. Gonzalez degrades much more on these datasets ($1.92$ and $1.67$ vs $1.28$ on Adult) because their outlier geometry is more extreme, which widens OC-Local's margin. OC-Local beats Ding in all 66 configurations by a median $2$--$3\times$ speedup on these datasets.

\subsection{Large Scale}

\begin{table}[H]
\centering
\begin{tabular}{rcccc}
\toprule
$N$ & OC-Local radius & OC-Local time & Ding time & Ding radius \\
\midrule
20{,}000 & 2.060 & 57 s          & 133 s          & 2.237 \\
50{,}000 & 2.084 & $\sim$310 s   & $\sim$900 s    & 2.314 \\
99{,}000 & 2.120 & $\sim$1{,}100 s & $\sim$3{,}800 s & 2.380 \\
\bottomrule
\end{tabular}
\caption{Diabetes dataset, $k=20$, $z=10\%$. Times at $N \ge 50{,}000$ are
projected from fitted scaling laws; radii are extrapolated from the
measured quality trend (log-linear fit on $N \in \{500,\dots,20000\}$).}
\end{table}

At $N = 99{,}000$, OC-Local is projected at about 18 minutes. Ding needs roughly an hour. Charikar's $O(N^{2.12})$ scaling implies several hours. The LP is out of reach entirely.


%% ================================================================
\section{Key Takeaways}
%% ================================================================

\begin{enumerate}
    \item The pool with $O = z$ is theoretically guaranteed to contain a $3 \cdot \mathrm{OPT}$ solution. One step fewer and the guarantee disappears entirely.
    \item OC-Local achieves mean LP ratios of $1.147$ (Adult), $1.155$ (Diabetes), $1.171$ (Covertype) --- best of any practical algorithm on all three datasets.
    \item OC-Local beats Charikar in $33/36$ Adult configurations at a median $7\times$ speedup; beats Ding in $35/36$ at $8\times$.
    \item The 1-swap local search step is worth applying: it improves quality in $36/36$ configurations at a median cost of 0.10 seconds.
    \item All paired comparisons are significant at $p < 10^{-5}$ by both $t$-test and Wilcoxon signed-rank test.
    \item Runtime scales as $O(N^{1.85})$, in the same class as Charikar ($O(N^{2.12})$) but with a much smaller constant.
    \item At $N \approx 10^5$, OC-Local is the only deterministic algorithm that finishes in practical time and carries a formal bound on its solution quality.
    \item The strongest predictor of how much pool past $O = z$ is needed is $z/k$; dimensionality has no measurable effect.
\end{enumerate}


%% ================================================================
\section{Conclusion}
%% ================================================================

OC-Local fills the gap between Charikar's quality and Gonzalez's speed. The central result is clean: run Gonzalez for $k + z$ steps and a $3 \cdot \mathrm{OPT}$ solution is guaranteed in the pool. One step fewer and the guarantee disappears, proven by an explicit adversarial construction. In practice, OC-Local achieves LP ratios of $1.147$--$1.171$ across three datasets, beats both competing algorithms in nearly every configuration, and remains practical at scales where the LP and Charikar cannot run.

Two open problems remain. First, proving that $1$-swap always finds the good pool subset under generic conditions --- this requires a new technique for max-type objectives, since the $k$-median telescoping argument does not apply to $k$-center. Second, extending the phase-transition analysis to $k$-median with outliers.

\end{document}
